%% ============================================================
%% AULA 9 — Atenção: quando a rede aprende onde olhar
%% GBC073 — Inteligência Computacional (FACOM/UFU)
%%
%% Esqueleto com esboço de conteúdo — a desenvolver.
%% Ementa (ficha SEI 5116656): item 2 — deep learning: introdução.
%%
%% Compilar:  latexmk -xelatex aula09.tex   (duas passadas!)
%% ============================================================
\documentclass[aspectratio=169, 11pt]{beamer}

\makeatletter
\def\input@path{{./}{../}}
\makeatother

\usetheme{InteligenciaComputacional}   % ANTES do babel: fontspec vem primeiro
\usepackage{babel}
\babelprovide[import, main]{brazilian}

\title[Aula 9 — Atenção]{Aula 9: Atenção — quando a rede aprende onde olhar}
\subtitle[GBC073 — Inteligência Computacional]{Self-attention \textbullet\ transformers \textbullet\ LLMs \textbullet\ o custo da escala}
\author[Prof.\ Marcelo Albertini]{Prof.\ Marcelo Keese Albertini}
\institute{GBC073 — Inteligência Computacional \textbullet\ Universidade Federal de Uberlândia}
\date{2026}

\begin{document}

% ------------------------------------------------------------ capa
\begin{frame}[plain, noframenumbering]
  \titlepage
\end{frame}

% ------------------------------------------------------------ roteiro
\begin{frame}{Roteiro da aula}
  \begin{itemize}\setlength{\itemsep}{0.5ex}
    \item<1-> \textbf{Atenção como média ponderada aprendida}
    \item<2-> \textbf{Self-attention: cada token olha todos os outros}
    \item<3-> \textbf{O transformer}
    \item<4-> \textbf{LLMs e o custo da escala}
  \end{itemize}
\end{frame}

% ------------------------------------------------------------
\section{Atenção como média ponderada aprendida}

\begin{frame}{Atenção como média ponderada aprendida}
  \begin{itemize}\setlength{\itemsep}{0.4ex}
    \item Na RNN, o passado inteiro precisa caber num único vetor de estado —
          gargalo.
    \item Atenção: a saída olha diretamente para todas as posições e decide
          o peso de cada uma.
    \item Consulta (\emph{query}), chave (\emph{key}) e valor (\emph{value}):
          a tríade da atenção.
    \item Os pesos de atenção são aprendidos — a rede aprende \emph{onde olhar}.
  \end{itemize}
\end{frame}

% ------------------------------------------------------------
\section{Self-attention}

\begin{frame}{Self-attention: cada token olha todos os outros}
  \begin{itemize}\setlength{\itemsep}{0.4ex}
    \item Uma sequência olha para si mesma: cada token pondera todos os
          demais, incluindo ele próprio.
    \item Caminho de informação de comprimento constante entre qualquer par
          de posições.
    \item Multi-cabeças: várias relações em paralelo (sintaxe, gênero,
          co-referência\ldots).
    \item Codificação posicional: sem recorrência, a ordem precisa ser
          injetada.
  \end{itemize}
\end{frame}

% ------------------------------------------------------------
\section{O transformer}

\begin{frame}{O transformer}
  \begin{itemize}\setlength{\itemsep}{0.4ex}
    \item Blocos repetidos: atenção + rede densa, com conexões residuais e
          normalização.
    \item Codificador--decodificador: tradução e afins.
    \item Só-decodificador: os modelos de linguagem modernos.
    \item Paralelismo total no treino — é a arquitetura que escala.
  \end{itemize}
\end{frame}

% ------------------------------------------------------------
\section{LLMs e o custo da escala}

\begin{frame}{LLMs e o custo da escala}
  \begin{itemize}\setlength{\itemsep}{0.4ex}
    \item Leis de escala: capacidade emerge com dados e computação juntos.
    \item Pré-treino em corpora imensos + ajuste fino para tarefas.
    \item O custo da escala: energia, dinheiro e a concentração de quem pode
          treinar.
    \item A revanche da aula 1: o perceptron, 65 anos depois.
  \end{itemize}
  \begin{quizbox}
    O custo computacional do self-attention cresce linearmente com o
    comprimento da sequência.

    \smallskip
    \opcoes{Verdadeiro \;/\; Falso}
  \end{quizbox}
\end{frame}

\end{document}
